Thermal runaway reactions can produce substantial quantities of gas.
LiionTR currently represents gas generation using a reduced-order
empirical formulation in which prescribed species yields are associated
with reaction progress.

This formulation provides the coupling required to study gas inventory,
internal pressure, and venting while keeping the gas-generation model
independent from the reaction kinetics and numerical solver.


\subsection{Gas species}

Each gas species $s$ is characterized by a molar mass

\begin{equation}
    M_s,
\end{equation}

expressed in \si{\kilogram\per\mole}.

The instantaneous gas state is represented by the mole inventory

\begin{equation}
    n_s,
\end{equation}

for each species.


\subsection{Gas inventory}

For a mixture containing $N_s$ species, the total gas amount is

\begin{equation}
    n_{\mathrm{tot}}
    =
    \sum_{s=1}^{N_s}
    n_s.
    \label{eq:gas_total_moles}
\end{equation}

The total gas mass is

\begin{equation}
    m_{\mathrm{gas}}
    =
    \sum_{s=1}^{N_s}
    n_s M_s.
    \label{eq:gas_total_mass}
\end{equation}

The mole fraction of species $s$ is

\begin{equation}
    x_s
    =
    \frac{
        n_s
    }{
        n_{\mathrm{tot}}
    },
    \label{eq:gas_mole_fraction}
\end{equation}

such that

\begin{equation}
    \sum_{s=1}^{N_s}
    x_s
    =
    1
\end{equation}

for a nonempty gas inventory.


\subsection{Mean molar mass}

The current venting model uses the mole-weighted mean molar mass

\begin{equation}
    \overline{M}
    =
    \sum_{s=1}^{N_s}
    x_s M_s.
    \label{eq:mean_molar_mass}
\end{equation}

Equivalently,

\begin{equation}
    \overline{M}
    =
    \frac{
        \displaystyle
        \sum_{s=1}^{N_s}
        n_s M_s
    }{
        \displaystyle
        \sum_{s=1}^{N_s}
        n_s
    }.
\end{equation}

This effective property provides the connection between the
species-resolved gas inventory and the single-mixture compressible-flow
model.


\subsection{Empirical reaction gas yields}

The current LiionTR gas-generation model assigns prescribed gas yields to
individual thermal runaway reactions.

For reaction $i$ and species $s$, the gas yield is denoted

\begin{equation}
    Y_{i,s},
\end{equation}

with units of

\begin{equation}
    \si{\mole\per\kilogram}
\end{equation}

of reacted material.

If reaction $i$ occupies an effective mass fraction $w_i$ of a cell with
mass $m_{\mathrm{cell}}$, the corresponding reacted-material rate is

\begin{equation}
    \frac{dm_{i,\mathrm{reacted}}}{dt}
    =
    m_{\mathrm{cell}}
    w_i
    \frac{d\alpha_i}{dt}.
    \label{eq:reacted_material_rate}
\end{equation}

The generation rate of species $s$ from reaction $i$ is therefore

\begin{equation}
    \dot{n}_{i,s}
    =
    Y_{i,s}
    m_{\mathrm{cell}}
    w_i
    \dot{\alpha}_i.
    \label{eq:species_generation_single_reaction}
\end{equation}


\subsection{Total species generation rate}

When multiple reactions generate the same gas species, their
contributions are summed.

The total generation rate of species $s$ is

\begin{equation}
    \dot{n}_{s,\mathrm{gen}}
    =
    \sum_{i=1}^{N_r}
    Y_{i,s}
    m_{\mathrm{cell}}
    w_i
    \dot{\alpha}_i,
    \label{eq:species_generation_total}
\end{equation}

where $N_r$ is the number of thermal runaway reactions.


\subsection{Cumulative gas generation}

If reaction $i$ evolves from an initial conversion
$\alpha_{i,0}$ to $\alpha_i$, the corresponding conversion change is

\begin{equation}
    \Delta \alpha_i
    =
    \alpha_i-\alpha_{i,0}.
\end{equation}

The cumulative gas amount associated with species $s$ is then

\begin{equation}
    n_{s,\mathrm{generated}}
    =
    \sum_{i=1}^{N_r}
    Y_{i,s}
    m_{\mathrm{cell}}
    w_i
    \left(
        \alpha_i-\alpha_{i,0}
    \right).
    \label{eq:cumulative_gas_generation}
\end{equation}

This expression is useful for closed-system calculations and for
verification of the integrated gas-generation rates.


\subsection{Closed-cell gas balance}

Before vent opening, all generated gas remains inside the modeled free
volume.

The species balance is

\begin{equation}
    \frac{dn_s}{dt}
    =
    \dot{n}_{s,\mathrm{gen}}.
    \label{eq:closed_gas_balance}
\end{equation}

The total mole inventory therefore increases according to

\begin{equation}
    \frac{dn_{\mathrm{tot}}}{dt}
    =
    \sum_s
    \dot{n}_{s,\mathrm{gen}}.
\end{equation}


\subsection{Gas balance after vent opening}

After vent opening, gas generation and gas discharge occur
simultaneously.

The species balance becomes

\begin{equation}
    \frac{dn_s}{dt}
    =
    \dot{n}_{s,\mathrm{gen}}
    -
    \dot{n}_{s,\mathrm{vent}}.
    \label{eq:open_gas_balance}
\end{equation}

The instantaneous gas composition may therefore evolve because of both
reaction-driven generation and removal through the vent.


\subsection{Coupling to pressure}

The gas inventory supplies the total gas amount used by the pressure
model.

For the current ideal-gas formulation,

\begin{equation}
    P
    =
    \frac{
        n_{\mathrm{tot}} R T
    }{
        V_{\mathrm{free}}
    }.
    \label{eq:gases_pressure_coupling}
\end{equation}

Thus, both gas generation and temperature evolution contribute directly
to cell pressurization.


\subsection{Coupling to venting}

The gas inventory also provides the mixture composition required by the
vent-flow model.

The mean molar mass is calculated using
\cref{eq:mean_molar_mass}, and the total molar discharge rate is
distributed between species according to their instantaneous mole
fractions.

This coupling is represented schematically in
\cref{fig:gas_architecture}.

\begin{figure}[htbp]
    \centering

    \begin{tikzpicture}[
        node distance=1.5cm and 2.3cm,
        block/.style={
            rectangle,
            rounded corners,
            draw,
            align=center,
            minimum width=3.2cm,
            minimum height=0.95cm
        },
        line/.style={
            -{Latex[length=3mm]},
            thick
        }
    ]

    \node[block] (reaction)
        {Reaction progress\\$\dot{\boldsymbol{\alpha}}$};

    \node[block, below=of reaction] (generation)
        {Gas generation\\$\dot{n}_{s,\mathrm{gen}}$};

    \node[block, below=of generation] (inventory)
        {Gas inventory\\$n_s$};

    \node[block, below left=of inventory] (pressure)
        {Pressure\\$P$};

    \node[block, below right=of inventory] (properties)
        {Mixture properties\\$x_s,\overline{M}$};

    \node[block, below=of $(pressure)!0.5!(properties)$] (vent)
        {Vent flow\\$\dot{n}_{s,\mathrm{vent}}$};

    \draw[line] (reaction) -- (generation);
    \draw[line] (generation) -- (inventory);
    \draw[line] (inventory) -- (pressure);
    \draw[line] (inventory) -- (properties);

    \draw[line] (pressure) -- (vent);
    \draw[line] (properties) -- (vent);

    \draw[line]
        (vent.east)
        -- ++(1.0,0)
        |- (inventory.east);

    \end{tikzpicture}

    \caption{
        Coupling between reaction-driven gas generation, the gas
        inventory, pressure, mixture properties, and vent discharge in
        the current LiionTR formulation.
    }
    \label{fig:gas_architecture}
\end{figure}


\subsection{Current assumptions}

The empirical gas-generation formulation assumes that:

\begin{itemize}
    \item gas yields are prescribed directly for individual reactions;
    \item species yields remain constant during each reaction;
    \item gas yields do not explicitly depend on temperature or pressure;
    \item gas-phase secondary chemistry is not represented;
    \item the gas mixture is spatially homogeneous;
    \item all gas species share a common temperature;
    \item the gas temperature is equal to the lumped cell temperature;
    \item condensation and phase separation are neglected;
    \item gas dissolution into condensed phases is neglected.
\end{itemize}


\subsection{Limitations of fixed species yields}

Thermal runaway gas composition can depend strongly on battery chemistry,
electrolyte composition, state of charge, temperature, pressure, and
secondary reactions.

Actual vent gases may contain, among other species,

\begin{equation}
    \mathrm{
        H_2,\,
        H_2O,\,
        CO,\,
        CO_2,\,
        CH_4,\,
        C_2H_4
    }.
\end{equation}

Additional fluorinated or phosphorus-containing species may also be
important depending on electrolyte chemistry.

A fixed species-yield model cannot generally predict these composition
changes from first principles.


\subsection{Future thermochemical formulation}

The intended long-term thermochemical architecture is to replace or
complement direct species yields with elemental source inventories.

Conceptually,

\begin{equation}
    \text{reaction progress}
    \longrightarrow
    \text{element inventory}
    \longrightarrow
    \text{chemical equilibrium}
    \longrightarrow
    \text{gas composition}.
\end{equation}

Relevant elemental inventories may include

\begin{equation}
    \mathrm{
        C,\,
        H,\,
        O,\,
        N,\,
        F,\,
        P,\,
        Li
    }.
\end{equation}

A thermochemical backend such as Cantera could then provide equilibrium
composition and mixture properties.

The present empirical implementation is retained because it provides a
simple, deterministic, and independently testable formulation while the
thermochemical architecture is developed.